\section{Identification and Estimation}\label{sec:identification}

Three pieces need identification: the type-specific cost distributions $F_c^k$, the auction-level UH component, and the entry counts. The last of these is directly observed---I compute Pre and Post means of bidders per auction by type and stratum. The first two require structure.

\subsection{Cost distributions $F_c^k$}

Under the English-reverse interpretation of Preg\~ao (Assumption~\ref{a:ipv}), the drop-out price of each loser point-identifies that bidder's cost. I therefore estimate $F_c^k$ by the empirical distribution of the losers' normalized drop-outs in each stratum (period $\times$ pharma $\times$ type). The losers-only estimator is unbiased as $N \to \infty$ but systematically omits the lowest cost draw per auction---the winner. In the baseline specification I therefore use the all-bidders ECDF (losers plus winner at her final clock price), which in a descending-clock auction approximates the winner's cost by the second-lowest exit $c_{(2)}$. This biases $F_c^k$ upward slightly (the winner's bid is above her cost), but less than losers-only biases it upward (which omits the lowest tail entirely). The robustness section reports a Turnbull NPMLE that treats the winner as left-censored at $c_{(2)}$ and confirms that the all-bidders approximation is tight.

A cross-modality consistency check supports the descending-clock specification. Under the same sample, Convite auctions are first-price sealed-bid; the \citet{gpv2000} inversion (GPV) applied to Convite bids recovers $F_c^k$ from the bid distribution. The two estimates---Convite GPV and Preg\~ao drop-out---coincide to within two percentage points at the median in pharma non-SME Pre, the stratum where the identifying assumption of cost-primitive stability holds most cleanly. In other strata the gap is wider and provides the empirical motivation for the UH deconvolution step.

\subsection{Auction-level unobserved heterogeneity}

The residual gap across modalities---up to twenty-three percentage points at the median in non-pharma---signals that the reference-price normalization does not fully absorb auction-level variation. Following \citet{krasnokutskaya2011}, I decompose $y_{it} = \log(b_{it}/p_t^{\mathrm{ref}})$ as $a_t + e_{it}$, where $a_t$ is common across bidders within auction $t$ and $e_{it}$ is bidder-specific (Assumption~\ref{a:uh}). A method-of-moments decomposition by stratum estimates $\hat\sigma_a^{2,k}$ and $\hat\sigma_e^{2,k}$; the intraclass correlation $\rho^k = \sigma_a^{2,k} / (\sigma_a^{2,k} + \sigma_e^{2,k})$ ranges from 0.47 to 0.62 across strata, in line with the 0.5--0.7 range that \citet{krasnokutskaya2011} reports for California highway procurement. A best-linear-predictor (BLP) shrinkage then extracts $\hat a_t = \frac{\hat\sigma_a^2 N_t}{\hat\sigma_a^2 N_t + \hat\sigma_e^2} \cdot (\bar y_t - \bar y)$ and delivers cleaned residuals $\hat e_{it} = y_{it} - \hat a_t$. The UH-clean bids $c_{\text{norm}}^{\text{clean}}$ are the inputs to all downstream structural estimation. After the correction, the cross-modality median gap closes by ninety-three percent in pharma SME Post and by twenty-three percent in pharma non-SME Post; in non-pharma the residual gap remains near the raw level, consistent with richer item-level heterogeneity that the variance decomposition captures only partially.

The BLP step assumes $a_t$ and $e_{it}$ are Gaussian, which is Krasnokutskaya's semi-parametric shortcut relative to the fully non-parametric Kotlarski (1967) deconvolution. The latter is implemented in the robustness section as a Turnbull nonparametric MLE of $F_c$ with winners treated as left-censored at $c_{(2)}$; the Turnbull CDF lies between the losers-only and all-bidders ECDFs in 24 of 24 comparisons across strata and quantiles (Table~\ref{tab:v3_turnbull_fc}), and delivers quantitatively indistinguishable BNE-decomposition numbers (Table~\ref{tab:v3_sensitivity_fc}).

\subsection{Entry}

Entry is not fully structural here. I take the observed average counts of SME and non-SME bidders per auction by (period $\times$ pharma) as the equilibrium arrival rates, model arrivals as Poisson with the corresponding rate, and compute the counterfactual BNE prices under the observed Post-period SME pool. The gap to a structural \citet{atheyseira2011} entry model with estimated entry costs is reported separately: a zero-profit condition applied to the simulated distribution of winning-margin profits delivers an implied entry cost per firm of R\$0.55 (SME, non-pharma) to R\$2.46 (non-SME, non-pharma), in the range of overhead costs for responding to a BEC tender. The details are in Table~\ref{tab:v3_entry_cost} and the s4 memo in the replication archive.

\subsection{Counterfactual simulation}

With $F_c^k$ and entry counts in hand, I simulate counterfactual prices Monte-Carlo-style. For each of $B$ auctions: draw $n^{\text{SME}} \sim \mathrm{Poisson}(\lambda^{\text{SME}})$ and analogously for non-SMEs, draw costs from $F_c^k$, sort, record $c_{(1)}$ (the winner's cost), $c_{(2)}$ (the realized price under revenue equivalence), and the winner's type. Aggregating over draws delivers means, distributions, and---for the welfare analysis---the allocative wedge $c_{(1)}^{S_3} - c_{(1)}^{S_1}$. Standard errors are computed via cluster bootstrap at the auction level with $B_{bs} = 500$ replicates (Section~\ref{sec:robustness}).
